\subsection{Contributions}

This paper makes three primary contributions to the CHI community:

\begin{enumerate}
\item \textbf{Empirical evidence of norm emergence patterns in LLM multi-agent systems.} We present a fully crossed $2\times2\times2\times2$ factorial experiment ($N=480$ runs) systematically varying communication structure, feedback loops, context architecture, and interaction density. We demonstrate that communication structure is the dominant determinant of norm emergence ($\eta^2=0.31$, $p<0.001$), and that model upgrades alone do not produce emergence.

\item \textbf{A typology of norm emergence trajectories.} We identify four distinct trajectories—Rapid Crystallization, Slow Convergence, Oscillatory Divergence, and Failure to Emerge—each associated with specific environmental configurations. This typology enables designers to predict emergence outcomes from platform configuration choices.

\item \textbf{The Norm Emergence Environmental Model (NEEM).} We propose NEEM as a design framework for the CHI community: a predictive model that maps environmental configuration to expected norm emergence outcomes. NEEM's central design principle is that \emph{communication infrastructure—guaranteed message delivery, bounded latency, shared context windows—produces norm emergence at low cost}, reframing the platform design problem as an infrastructure investment rather than a model capability problem.
\end{enumerate}

We further validate NEEM against the Tsinghua Moltbook 93\% non-response baseline, confirming ecological validity. Our findings have direct implications for the design of multi-agent HCI systems: copilots, collaborative AI assistants, autonomous negotiation platforms, and human-AI teams.